% Unit 6 exam -- 20 MCQ + 2 short FRQ = 28 pts, 48 min.
\documentclass{apchem-exam}

\apcunit{6}
\apcunittitle{Thermochemistry}
\apcdoctitle{Unit Exam}
\apcsubtitle{CED 6.1--6.9 \textbullet\ Fri Dec 11 \textbullet\ 48 minutes \textbullet\ 28 points}

\begin{document}
\apctitleblock

\begin{apcnote}[Format]
Section I: 20 multiple choice, \textbf{no calculator}, 30 minutes.
Section II: two short free response, calculator permitted, 18 minutes.
Constants are provided with Section II.
\end{apcnote}

\examsection{I}{Multiple Choice}{20 questions \textbullet\ 30 minutes \textbullet\ 1 point each}{\nocalculator}

\begin{mcq}
  A reaction is run in water and the water's temperature rises. The
  reaction is
  \choices
    \correct{exothermic, with $\Delta H < 0$}
    \choice{endothermic, with $\Delta H > 0$}
    \choice{exothermic, with $\Delta H > 0$}
    \choice{at thermal equilibrium}
\end{mcq}
\why{The water is the surroundings; it warmed because the system released
energy.}
\distractor{B}{the most common error in the unit --- reading the
thermometer as if it were in the system}

\begin{mcq}
  For an endothermic process, $q_{\text{system}}$ is
  \choices
    \correct{positive, and the surroundings cool}
    \choice{positive, and the surroundings warm}
    \choice{negative, and the surroundings cool}
    \choice{negative, and the surroundings warm}
\end{mcq}
\why{Energy enters the system, so it must leave the surroundings.}

\begin{mcq}
  On an energy diagram, products drawn above reactants indicates
  \choices
    \choice{an exothermic reaction}
    \correct{an endothermic reaction}
    \choice{a fast reaction}
    \choice{a slow reaction}
\end{mcq}
\why{The products hold more enthalpy, so energy was absorbed.}

\begin{mcq}
  The height of the barrier on an energy diagram affects
  \choices
    \correct{the rate but not $\Delta H$}
    \choice{$\Delta H$ but not the rate}
    \choice{both the rate and $\Delta H$}
    \choice{neither}
\end{mcq}
\why{$\Delta H$ is fixed by the reactant and product levels alone.}

\begin{mcq}
  Two objects at different temperatures are placed in contact. Energy
  flows
  \choices
    \correct{from the hotter to the cooler until both reach the same
      temperature}
    \choice{from the cooler to the hotter}
    \choice{in whichever direction the larger object dictates}
    \choice{only if they are the same material}
\end{mcq}
\why{Always hot to cold, ending at thermal equilibrium.}

\begin{mcq}
  At thermal equilibrium,
  \choices
    \correct{energy is still exchanged, but at equal rates in both
      directions}
    \choice{all molecular motion stops}
    \choice{energy exchange ceases entirely}
    \choice{the two objects must have equal internal energy}
\end{mcq}
\why{Equilibrium is dynamic --- no \emph{net} transfer, not no transfer.}

\begin{mcq}
  In the expression $q = mc\Delta T$ applied to a coffee-cup calorimeter,
  $m$ and $c$ refer to
  \choices
    \correct{the water}
    \choice{the reactants}
    \choice{the calorimeter cup}
    \choice{the limiting reactant only}
\end{mcq}
\why{The water is the surroundings whose temperature was measured.}

\begin{mcq}
  Water's specific heat is unusually large mainly because of
  \choices
    \correct{hydrogen bonding}
    \choice{its low molar mass}
    \choice{its bent shape}
    \choice{its high boiling point}
\end{mcq}
\why{Energy goes into disrupting the hydrogen-bond network.}

\begin{mcq}
  Heating \SI{50.0}{\gram} of water by \SI{60.0}{\celsius} requires about
  \choices
    \choice{\SI{1.25}{\kilo\joule}}
    \correct{\SI{12.5}{\kilo\joule}}
    \choice{\SI{125}{\kilo\joule}}
    \choice{\SI{250}{\kilo\joule}}
\end{mcq}
\why{$(50.0)(4.18)(60.0) \approx 12{,}500$~J.}

\begin{mcq}
  A calorimetry experiment gives $q_{\text{water}} = +2170$~J with
  \SI{0.0250}{\mole} reacting. Then $\Delta H$ is about
  \choices
    \choice{$+86.9$ kJ/mol}
    \correct{$-86.9$ kJ/mol}
    \choice{$-2.17$ kJ/mol}
    \choice{$-54.3$ kJ/mol}
\end{mcq}
\why{Reverse the sign, then divide by moles.}
\distractor{A}{forgot the sign reversal from water to reaction}

\begin{mcq}
  Which expression gives the energy to melt a sample at its melting point?
  \choices
    \choice{$mc\Delta T$}
    \correct{$n\Delta H_{\text{fus}}$}
    \choice{$n\Delta H_{\text{vap}}$}
    \choice{$mc\Delta T + n\Delta H_{\text{fus}}$}
\end{mcq}
\why{At the melting point the temperature is not changing.}

\begin{mcq}
  During boiling, the temperature stays constant because the added energy
  \choices
    \correct{overcomes intermolecular forces rather than increasing kinetic
      energy}
    \choice{is lost to the surroundings}
    \choice{breaks covalent bonds within the molecules}
    \choice{is converted into pressure}
\end{mcq}
\why{Temperature tracks average kinetic energy, which is unchanged.}
\distractor{C}{boiling separates molecules; it does not break O--H bonds}

\begin{mcq}
  For water, $\Delta H_{\text{vap}}$ compared with
  $\Delta H_{\text{fus}}$ is
  \choices
    \correct{much larger, because vaporizing separates the molecules
      entirely}
    \choice{much smaller}
    \choice{about equal}
    \choice{larger only above \SI{100}{\celsius}}
\end{mcq}
\why{\SI{40.7}{} against \SI{6.02}{\kilo\joule\per\mole}.}

\begin{mcq}
  For $\ce{2CO + O2 -> 2CO2}$, $\Delta H = -566$~kJ. For
  $\ce{2CO2 -> 2CO + O2}$, $\Delta H$ is
  \choices
    \choice{$-566$ kJ}
    \correct{$+566$ kJ}
    \choice{$-283$ kJ}
    \choice{$+283$ kJ}
\end{mcq}
\why{Reversing a reaction flips the sign.}

\begin{mcq}
  Doubling every coefficient in a thermochemical equation
  \choices
    \correct{doubles $\Delta H$}
    \choice{leaves $\Delta H$ unchanged}
    \choice{halves $\Delta H$}
    \choice{changes the sign of $\Delta H$}
\end{mcq}
\why{$\Delta H$ is extensive.}

\begin{mcq}
  Hess's law is valid because enthalpy is
  \choices
    \correct{a state function}
    \choice{always negative}
    \choice{independent of temperature}
    \choice{measured in kilojoules}
\end{mcq}
\why{Path-independence is exactly what a state function means.}

\begin{mcq}
  The enthalpy of formation of \ce{O2(g)} at standard conditions is
  \choices
    \correct{zero, by definition}
    \choice{$-286$ kJ/mol}
    \choice{$+495$ kJ/mol}
    \choice{not defined}
\end{mcq}
\why{An element in its standard state needs no forming.}

\begin{mcq}
  $\Delta H$ from bond enthalpies is calculated as
  \choices
    \correct{bonds broken minus bonds formed}
    \choice{bonds formed minus bonds broken}
    \choice{products minus reactants}
    \choice{the sum of all bond energies}
\end{mcq}
\why{Breaking costs energy, so those terms enter positively.}
\distractor{B}{the reversed convention is the standard trap here}

\begin{mcq}
  A bond-enthalpy estimate and a formation-enthalpy calculation for the
  same reaction differ by about 20~kJ. The best explanation is that
  \choices
    \correct{bond enthalpies are averages that ignore molecular environment}
    \choice{one of the calculations must contain an arithmetic error}
    \choice{the reaction is not at standard conditions}
    \choice{bond enthalpies apply only to gases}
\end{mcq}
\why{Averaging is exactly why the two routes disagree slightly.}

\begin{mcq}
  For $\ce{H2 + Cl2 -> 2HCl}$ with $D$(H--H) $=432$,
  $D$(Cl--Cl) $=239$, $D$(H--Cl) $=427$~kJ/mol, $\Delta H$ is
  \choices
    \correct{$-183$ kJ}
    \choice{$+183$ kJ}
    \choice{$-244$ kJ}
    \choice{$+854$ kJ}
\end{mcq}
\why{$671 - 854$.}

\examsection{II}{Free Response}{2 questions \textbullet\ 18 minutes}{\calculatorok}

\begin{apcnote}[Data]
$c$(water) $= \SI{4.18}{\joule\per\gram\per\celsius}$ \textbullet\
$\Delta H_{\text{fus}} = \SI{6.02}{\kilo\joule\per\mole}$ \textbullet\
$\Delta H_{\text{vap}} = \SI{40.7}{\kilo\joule\per\mole}$ \textbullet\
$M(\ce{H2O}) = \SI{18.02}{\gram\per\mole}$. \\
$\Delta H_f^\circ$ (kJ/mol): \ce{CO2(g)} $-393.5$ \textbullet\
\ce{H2O(l)} $-286$ \textbullet\ \ce{CH4(g)} $-75$ \textbullet\
\ce{O2(g)} $0$.
\end{apcnote}

% ---- FRQ 1 -----------------------------------------------------------------
\frq{short}{4}{CED 6.4, 6.6 \textbullet\ calorimetry}

A \SI{0.0400}{\mole} sample of methane is completely burned in a
calorimeter containing \SI{500.0}{\gram} of water. The water's temperature
rises from \SI{21.00}{\celsius} to \SI{38.05}{\celsius}.

\begin{parts}
  \item Calculate the heat absorbed by the water. \workspace[2cm]
        \keyonly{\begin{apcanswer}$\Delta T = 17.05$;
        $q = (500.0)(4.18)(17.05) = \mathbf{35{,}600}~\text{J}
        = 35.6~\text{kJ}$\end{apcanswer}}
  \item Calculate $\Delta H$ of combustion in \si{\kilo\joule\per\mole}.
        \workspace[2.2cm]
        \keyonly{\begin{apcanswer}$q_{\text{rxn}} = -35.6$~kJ;
        $\Delta H = -35.6/0.0400 =
        \mathbf{-891}~\si{\kilo\joule\per\mole}$\end{apcanswer}}
  \item The accepted value is $-890.5$~\si{\kilo\joule\per\mole}. Give one
        reason a real calorimetry experiment usually yields a magnitude
        that is too small. \answerlines[2]{Some energy is absorbed by the
        calorimeter itself or escapes to the surroundings rather than
        heating the water, so the measured temperature rise --- and the
        energy computed from it --- under-reports the reaction.}
\end{parts}

\begin{rubric}
  \rpoint{1}{(a) $\Delta T = 17.05$ and $q \approx 35.6$~kJ}
  \rpoint{2}{(b) 1 pt sign reversal to $q_{\text{rxn}}$; 1 pt division by
    moles with the answer in \si{\kilo\joule\per\mole} --- an answer in bare
    kJ earns 1}
  \rpoint{1}{(c) heat loss to the calorimeter or surroundings}
\end{rubric}

% ---- FRQ 2 -----------------------------------------------------------------
\frq{short}{4}{CED 6.5, 6.9 \textbullet\ phase changes and Hess's law}

\begin{parts}
  \item Calculate the energy needed to convert \SI{36.0}{\gram} of ice at
        \SI{0}{\celsius} into liquid water at \SI{25.0}{\celsius}.
        \workspace[2.8cm]
        \keyonly{\begin{apcanswer}Melt: $n = 36.0/18.02 = 2.00$~mol,
        $q = (2.00)(6.02) = 12.0$~kJ. Warm:
        $(36.0)(4.18)(25.0) = 3762$~J $= 3.76$~kJ.
        \textbf{Total} $= \mathbf{15.8}$~kJ\end{apcanswer}}
  \item Explain why the calculation needs two separate expressions.
        \answerlines[2]{Melting happens at constant temperature, so it uses
        $n\Delta H_{\text{fus}}$; warming the liquid changes temperature
        within a single phase, so it uses $mc\Delta T$. Neither expression
        covers both stages.}
  \item Given $\ce{C(gr) + O2 -> CO2}$, $\Delta H = -393.5$~kJ, and
        $\ce{CO + \tfrac{1}{2}O2 -> CO2}$, $\Delta H = -283.0$~kJ,
        find $\Delta H$ for $\ce{C(gr) + \tfrac{1}{2}O2 -> CO}$.
        \workspace[2.2cm]
        \keyonly{\begin{apcanswer}Reverse the second, flipping its sign to
        $+283.0$, and add to the first:
        $-393.5 + 283.0 = \mathbf{-110.5}~\text{kJ}$\end{apcanswer}}
\end{parts}

\begin{rubric}
  \rpoint{2}{(a) 1 pt melting term 12.0~kJ; 1 pt warming term 3.76~kJ and
    total 15.8~kJ --- using $mc\Delta T$ for the melting stage earns 0 for
    the first point}
  \rpoint{1}{(b) identifies constant temperature versus changing
    temperature as the reason}
  \rpoint{1}{(c) $-110.5$~kJ with the second equation reversed and
    re-signed}
\end{rubric}

\examtotal

\end{document}
